\section{Experiments}
\label{sec:experiments}

\subsection{Experimental Setup}
All experiments use the train/validation/test splits for segmentation and classification described in Section~\ref{sec:data}.

For evaluation, segmentation performance is measured using Dice coefficient, IoU, and the 95th percentile Hausdorff distance (HD95). Dice and IoU measure overlap between the predicted and ground-truth tumor masks. HD95 measures boundary distance by reporting the 95th percentile of surface-to-surface distances, making it less sensitive to extreme outlier pixels than the standard Hausdorff distance. Classification performance is measured using accuracy, precision, recall, and F1-score.

The hyperparameters in Table~\ref{tab:optimization_details} were selected using validation performance and computational constraints. For classification, the learning rate and batch size were selected using a lightweight grid search over learning rates $\{10^{-3}, 5 \times 10^{-4}, 10^{-4}, 5 \times 10^{-5}\}$ and batch sizes $\{16, 32, 64\}$, with validation F1 as the selection criterion. The remaining settings follow the training configuration used in the pipeline.

\begin{table}[H]
    \centering
    \small
    \caption{Training hyperparameters used in the main experiments.}
    \label{tab:optimization_details}
    \setlength{\tabcolsep}{3pt}
    \makebox[\linewidth][c]{%
    \begin{tabular}{lcc}
    \hline
    Hyperparameter & Segmentation & Classification \\
    \hline
    Input size & $256 \times 256$ & $224 \times 224$ \\
    Optimizer & Adam & Adam \\
    Learning rate & $3 \times 10^{-4}$ & $5 \times 10^{-5}$ \\
    Weight decay & $1 \times 10^{-5}$ & $1 \times 10^{-4}$ \\
    Batch size & 8 & 32 \\
    Maximum epochs & 130 & 50 \\
    Early-stopping patience & 12 epochs & 10 epochs \\
    Scheduler & ReduceLROnPlateau & ReduceLROnPlateau \\
    Selection metric & Validation Dice & Validation F1 \\
    \hline
    \end{tabular}%
    }
\end{table}

Early-stopping patience denotes the number of consecutive epochs without validation-loss improvement before training is stopped.




\subsection{Segmentation Comparison}
To select the segmentation backbone, we compare U-Net, ResUNet, and Attention U-Net under the same data split, training protocol, and evaluation metrics. Table~\ref{tab:seg-results} shows that ResUNet achieves the best Dice and IoU, while U-Net obtains the lowest HD95. Since Dice and IoU better reflect overall tumor-region overlap, we select ResUNet as the final localization backbone.

\begin{table}[H]
\centering
\caption{Segmentation results on the test set.}
\label{tab:seg-results}
\begin{tabular}{lccc}
\hline
Model & Dice $\uparrow$ & IoU $\uparrow$ & HD95 $\downarrow$ \\
\hline
U-Net & 0.7102 & 0.6230 & \textbf{11.55} \\
ResUNet & \textbf{0.7431} & \textbf{0.6558} & 11.85 \\
Attention U-Net & 0.7360 & 0.6472 & 14.56 \\
\hline
\end{tabular}
\end{table}

Qualitative inspection suggests that the models usually identify the main tumor region, but common failure cases include small tumors, low-contrast boundaries, and irregular lesions where part of the tumor is missed.




\subsection{Classification Ablation}
We evaluate several classification settings to understand the effect of segmentation-guided cropping. The full-image ResNet50 baseline removes segmentation guidance and classifies the original MRI image directly. The cropped multi-stage variants use the predicted tumor ROI, with different bounding-box margins, to test whether the crop size limits classification performance. Finally, the dual-input model combines both the cropped ROI and the original full image, allowing the classifier to use localized tumor information while retaining global anatomical context.

\begin{table}[H]
    \centering
    \small
    \caption{Classification ablation results on the test set.}
    \label{tab:cls-results}
    \setlength{\tabcolsep}{1pt}
    \makebox[\linewidth][c]{%
    \begin{tabular}{lcccc}
    \hline
    Input setting & Accuracy & Precision & Recall & F1 \\
    \hline
    Full-image ResNet50 baseline & \textbf{0.9915} & \textbf{0.9914} & \textbf{0.9909} & \textbf{0.9911} \\
    Cropped ROI, 10-pixel margin & 0.9516 & 0.9498 & 0.9502 & 0.9496 \\
    Cropped ROI, 15-pixel margin & 0.9579 & 0.9476 & 0.9532 & 0.9504 \\
    Cropped ROI, 20-pixel margin & 0.9621 & 0.9520 & 0.9674 & 0.9596 \\
    Dual input: cropped ROI + full image & 0.9900 & 0.9895 & 0.9895 & 0.9895 \\
    \hline
    \end{tabular}%
    }
    \end{table}

Table~\ref{tab:cls-results} shows that increasing the ROI margin improves the cropped-only classifier, with the 20-pixel margin giving the best cropped result. However, the improvement from 10 to 20 pixels is modest compared with the remaining gap between cropped-only classification and the full-image baseline. Supplementary Figures~\ref{fig:supp_cm_resunet_single} and~\ref{fig:supp_cm_resunet_dual} show that the cropped-only model makes more cross-class mistakes among tumor categories, while the dual-input model reduces these errors by restoring global image context. This suggests that the main limitation is not only the crop margin, but also the loss of global anatomical context and the propagation of segmentation errors into the classifier input.

The dual-input model substantially reduces this gap by combining the localized tumor crop with the original full image. Its test F1-score reaches 0.9895, much closer to the full-image baseline than any cropped-only variant. The training curves in Supplementary Figures~\ref{fig:supp_dual_training_curves} and~\ref{fig:supp_dual_metric_curves} show stable optimization, with validation metrics converging to a high level. Although the full-image baseline still achieves the highest numerical performance, we use the dual-input setting for the final explainable pipeline because it preserves the interpretability benefits of segmentation-guided localization while recovering most of the global context needed for accurate classification.





\subsection{VLM Summary Generation Comparison}
We compare four VLMs for the summary-generation component: BLIP-2, InstructBLIP, CLIP, and MedGemma. The exact model checkpoints are listed in Table~\ref{tab:vlm-results}. The generated summaries are evaluated using reference-based text-similarity metrics and task-specific summary-quality metrics.

ROUGE-L measures the longest common subsequence overlap between the generated summary and the reference summary, capturing similarity in wording and sentence structure. METEOR measures word-level alignment while also considering more flexible matches such as stemming and synonyms. The keyword-hit rate measures whether clinically relevant tumor terms appear in the generated summary. The composite score summarizes task-specific summary quality by combining keyword coverage, confidence reporting, and length appropriateness.

\begin{table}[H]
    \centering
    \small
    \caption{VLM summary-generation comparison.}
    \label{tab:vlm-results}
    \setlength{\tabcolsep}{2pt}
    \makebox[\linewidth][c]{%
    \begin{tabular}{lcccc}
    \hline
    Model & ROUGE-L & METEOR & Keyword hit & Composite \\
    \hline
    blip2-opt-2.7b & 0.1526 & 0.1264 & 0.4909 & 0.7425 \\
    instructblip-vicuna-7b & 0.2408 & 0.2023 & 0.4491 & 0.6991 \\
    clip-vit-base-patch32 & \textbf{0.2716} & 0.2358 & 0.4109 & 0.7025 \\
    medgemma-1.5-4b-it & 0.1694 & \textbf{0.3431} & \textbf{0.6600} & \textbf{0.8300} \\
    \hline
\end{tabular}%
}
    \end{table}

Table~\ref{tab:vlm-results} shows that \texttt{clip-vit-base-patch32} achieves the highest ROUGE-L score, indicating the strongest longest-subsequence overlap with the reference summaries. However, \texttt{medgemma-1.5-4b-it} achieves the best METEOR score, keyword-hit rate, and composite score. This suggests that MedGemma produces summaries with better clinically relevant terminology and stronger task-specific alignment, even if its surface-level wording is less similar to the reference summaries. Therefore, we use medgemma-1.5-4b-it as the summary-generation model in the final pipeline.





\subsection{Explainability Analysis}
Grad-CAM visualizations are used to inspect whether the classifier focuses on clinically meaningful regions. In successful examples, the heatmap overlaps with the predicted segmentation mask, suggesting that the segmentation and classification stages rely on similar tumor evidence. In failure cases, the heatmap may activate outside the predicted tumor area or on high-contrast background structures, indicating possible reliance on spurious cues.

These visual checks complement the quantitative ablations by showing not only whether the model is correct, but also where its evidence comes from. This is important for the final pipeline because the prediction, localized ROI, Grad-CAM attention, and generated summary should be consistent with one another.
Supplementary Figures~\ref{fig:supp_debug_notumor}, \ref{fig:supp_debug_pituitary}, \ref{fig:supp_real_notumor}, \ref{fig:supp_real_pituitary}, and~\ref{fig:supp_real_glioma} show qualitative examples from the final pipeline, including the original MRI, segmentation output, Grad-CAM visualization, and generated summary.




